% Source: Notion page “5. 语义增强的轻量级骨架行为识别模型设计”. Generated without rewriting its content.
\chapter{语义增强的轻量级骨架行为识别模型设计}

本系统面向校园场景中的细粒度行为识别任务，在轻量级骨架图卷积网络 ProtoGCN 的基础上，引入生成式动作描述驱动的图—文本语义对齐机制。一方面，利用 ProtoGCN 对人体关节之间的动态拓扑关系进行建模，通过运动拓扑增强和原型重构突出相似动作之间的细粒度差异；另一方面，借鉴 Generative Action Description Prompts（GAP）的跨模态对比学习思想，将行为类别的自然语言语义作为训练阶段的辅助监督，使骨架特征不仅具有运动结构上的判别能力，同时获得更明确的动作语义约束。

针对“嬉戏追逐—冲突追逐”“嬉戏推搡—冲突推搡”等骨架运动模式高度相似的校园行为，本系统最终形成“拓扑结构判别 + 原型差异强化 + 动作语义约束”的三级特征学习机制。

\section{模型训练总体思路}

本系统采用“大规模动作数据预训练 + 校园六分类全量微调”的两阶段训练策略，而不是直接利用规模较小的校园行为数据从零训练模型。

整体训练流程可以表示为：

\begin{NotionCode}[plain text]
阶段一：Kinetics-400 骨架动作预训练
Kinetics Skeleton 2D
        ↓
骨架采样 / Kinetics_Transform / coco_new 图构建
        ↓
ProtoGCN Backbone
GCN Block + MSTCN + MTE
        ↓
Prototype Reconstruction Network（PRN）
        ↓
400 类行为分类
        ↓
Cross Entropy + CSCL
        ↓
得到 K400 预训练 Backbone θ_K400


阶段二：Campus6 全参数微调
校园六类骨架数据
        ↓
相同的骨架预处理与 coco_new 图结构
        ↓
加载 θ_K400
        ↓
ProtoGCN + PRN
        ├────────────→ 六分类任务分支
        │
        └────────────→ GAP 身体部位语义监督分支
                         （仅训练阶段）
        ↓
CE + CSCL + GAP Semantic Loss
        ↓
六分类 Campus6 Student
\end{NotionCode}

这样的训练方式主要解决两个问题：一方面，Kinetics-400 提供大量通用人体动作样本，使模型首先学习稳定的关节运动模式、人体拓扑关系及时间动态；另一方面，校园数据微调进一步把这些通用运动知识适配到 \(l_w\)alk、\(l_r\)un、\(l_c\)hase、\(l_p\)ush、\(t_c\)hase、\(t_p\)ush 六个目标类别，尤其强化“嬉戏—冲突”等细粒度行为之间的区分能力。六类标签与当前工程定义保持一致。

\section{阶段一：基于 Kinetics-400 的骨架动作预训练}

\subsection{骨架输入与时空图构建}

预训练阶段使用 Kinetics Skeleton 2D 数据。

对于每个动作样本，首先对完整骨架序列进行均匀时间采样，将序列统一为 100 帧；对于多人场景，最多保留 2 名主要人物。随后通过 Kinetics Transform 将关键点统一转换至 ProtoGCN 使用的 \(o_n\)ew 图结构。

当前模型使用 joint modality，即直接以人体关键点位置作为基础骨架特征。Kinetics-400 配置中采用 \(p_l\)en=100、\(m_p\)erson=2 和 \(o_n\)ew 图结构。

对于连续骨架序列，可以将其抽象为时空图：

G=(V,E),G=(V,E),

其中：

\begin{itemize}
\item \(V\) 表示人体关键点节点；
\item 空间边描述同一帧中人体骨骼的天然连接关系；
\item 时间边描述连续时间帧中同一关键点之间的运动关联。
\end{itemize}

因此，原始人体骨架序列被转换为一个 Spatio-Temporal Graph，后续所有特征学习均建立在该图结构之上。


\section{ProtoGCN 主干：从关节运动到细粒度拓扑特征}

普通 GCN 能够利用人体骨骼连接传播关节信息，但对于 \(l_p\)ush 与 \(t_p\)ush、\(l_c\)hase 与 \(t_c\)hase 这类行为而言，参与运动的身体部位往往高度相似。

真正决定类别的并不只是“哪些关节发生了运动”，还包括：

\begin{itemize}
\item 关节运动幅度；
\item 不同身体部位之间的协同关系；
\item 人与人之间距离变化；
\item 局部关节运动的细微差异。
\end{itemize}

因此，本系统采用 ProtoGCN 作为骨架特征提取主干，通过 Motion Topology Enhancement（MTE）+ Prototype Reconstruction Network（PRN）+ Class-Specific Contrastive Learning（CSCL） 对动作关系进行进一步建模。

\subsection{GCN Block 与多尺度时间建模}

ProtoGCN 的基本网络单元由图卷积和多尺度时间卷积组成。

其中，GCN 主要负责建模同一时间帧内不同人体关节之间的空间依赖关系；MSTCN 则从多个时间尺度建模人体动作动态，使模型同时获取：

\[
\text{Spatial Relation}
+
\text{Temporal Motion}.
\]

即一方面判断“身体哪些部位存在运动关联”，另一方面学习“这些关联如何随时间发生变化”。

当前 ProtoGCN 实现也是按照 GCN、MSTCN 和残差连接组成多层时空特征提取网络。


\section{MTE：增强样本相关的细粒度运动拓扑}

传统 GCN 大多基于固定人体骨骼结构或跨样本共享的可学习图结构进行建模，但是不同动作样本中的真实关节关系并不完全相同。

例如，“正常行走”和“正常跑步”都涉及双腿交替运动，但跑步动作中的膝关节弯曲程度、步幅、上肢摆动幅度以及躯干移动速度明显更大。

因此 ProtoGCN 通过 Motion Topology Enhancement，进一步建立与当前样本相关的动态关节关系。

设某层输入特征为 \(\mathbf H\)，首先计算 Query 与 Key：

\[
H^{Q}
=
\operatorname{AvgPool}_{T}
(\mathbf H\mathbf W_Q),
\qquad
\mathbf H^{K}
=
\operatorname{AvgPool}_{T}
(\mathbf H\mathbf W_K).
\]

随后建模当前样本中的整体关节相关性：

\[
\mathbf A_{intra}
=
\phi_{\langle intra\rangle}
\left(
\mathbf H^{Q}
(\mathbf H^{K})^{\top}
\right).
\]

同时通过两两特征差异进一步描述关节之间的细粒度区别：

\[
\mathbf A_{inter}
=
\phi_{\langle inter\rangle}
\left(
\mathcal T_1(\mathbf H^Q)
-
\mathcal T_2(\mathbf H^K)
\right).
\]

最终将人体基础拓扑与动态学习到的拓扑共同用于图卷积：

\[
\mathbf H^{(l)}
=
\sigma
\left[
(
\mathbf A_0
+
\mathbf A_{intra}
+
\mathbf A_{inter}
)
\mathbf H^{(l-1)}
\mathbf W^{(l)}
\right].
\]

其中：

\[
\mathbf A_0
\]

描述人体共有的基础骨骼连接关系；

\[
\mathbf A_{intra}
\]

描述当前动作内部不同关节之间的整体相关性；

\[
\mathbf A_{inter}
\]

进一步保留不同关节运动模式之间的细粒度差异。

因此 MTE 的作用并不是简单地增加一个动态邻接矩阵，而是让模型同时学习“哪些关节共同运动”以及“它们具体如何不同地运动”。上述形式与 ProtoGCN 中的 MTE 定义保持一致。


\section{PRN：基于运动原型重构高层关节关系}

经过多层 GCN 与 MTE 后得到的高层动态图结构仍然可能包含姿态检测误差、遮挡以及偶然动作产生的无关关系。

因此 ProtoGCN 在高层动态图之后引入 Prototype Reconstruction Network（PRN）。

设最终学习到的高层关系图为：

\[
\mathbf A
\in
\mathbb R^{N\times N\times C}.
\]

将其展开得到：

\[
\mathbf X
\in
\mathbb R^{N^2\times C}.
\]

这里 \(\mathbf X\) 的每一行都可以理解为一个“关节对”的高维关系表示。

PRN 内部维护一个可学习的原型记忆空间。首先根据当前关系特征计算其对不同 prototype 的响应：

\[
R
=
\operatorname{softmax}
\left(
\mathbf X
\mathbf W_{query}^{\top}
\right)
\in
\mathbb R^{N^2\times n_{pro}}.
\]

随后按照这些响应权重重新组合 prototype：

\[
W_{memory}
\in
\mathbb R^{N^2\times C}.
\]

因此可以将整个过程概括为：

\[
\mathbf X
\rightarrow
\text{Prototype Assignment}
\rightarrow
\mathbf R
\rightarrow
\text{Prototype Memory}
\rightarrow
\mathbf Z.
\]

这里的 prototype 并不是某个固定动作类别，也不是简单对应“手”“脚”等身体区域，而是一种模型自动学习得到的抽象关节关系模式。

通过 prototype bottleneck，高层图关系不再完全依赖单一样本，而需要利用训练数据中学习到的典型运动关系重新表达，从而抑制部分偶然噪声，并突出更加稳定、具有判别性的运动模式。

当前模型使用：

\[
n_{pro}=400,
\]

即维护 400 个运动关系 prototype。


\section{CSCL：增强原型重构结果的类别判别能力}

PRN 保证了高层图关系可以通过一组 prototype 进行重构，但是仅完成 prototype 重构并不能保证这些关系一定具有类别判别能力。

因此 ProtoGCN 又引入 Class-Specific Contrastive Learning（CSCL）。

对于第 \(k\) 类动作，设其类别聚合中心为：

\[
\mathbf m_k.
\]

当前 batch 中属于第 \(k\) 类的平均特征记为：

\[
\bar{\mathbf f}_k.
\]

类别中心采用动量更新：

\[
\mathbf m_k
=
\alpha\mathbf m_k
+
(1-\alpha)
\bar{\mathbf f}_k.
\]

对于类别 kk 中的某一个样本 \(\mathbf f\)\_\{k,j\}，定义：

\[
\mathcal L_{CSC}
=
-\log
\frac{
\exp(
\mathbf f_{k,j}
\cdot
\mathbf m_k/\tau
)
}{
\exp(
\mathbf f_{k,j}
\cdot
\mathbf m_k/\tau
)
+
\sum_{l\neq k}
\exp(
\mathbf f_{k,j}
\cdot
\mathbf m_l/\tau
)
}.
\]

其本质包括两个方向：

Compactness：类内紧凑性

\[
f_{k,j}
\rightarrow
\mathbf m_k
\]

使相同行为类别的高层关系更加接近。

Dispersion：类间分离性

\[
f_{k,j}
\not\rightarrow
\mathbf m_l,\qquad l\neq k
\]

使不同动作类别的高层关系尽可能分离。

对于校园行为而言，这种机制特别适用于区分运动结构高度相似的类别，例如：

\[
\text{playful chase}
\leftrightarrow
\text{conflict chase}
\]

以及

\[
\text{playful push}
\leftrightarrow
\text{conflict push}.
\]

因此 Kinetics-400 预训练阶段的主要优化目标为：

\[
\boxed{
\mathcal L_{\text{pretrain}}
=
\mathcal L_{CE}^{K400}
+
\lambda_{CSC}
\mathcal L_{CSC}
}
\]

当前 ProtoGCN 中：

\[
C=0.2.\lambda_{CSC}=0.2.
\]

因此：

\[
\boxed{
\mathcal L_{\text{pretrain}}
=
\mathcal L_{CE}^{K400}
+
0.2
\mathcal L_{CSC}
}
\]

这也与仓库当前分类 Head 中的实际损失组合一致。


\section{阶段一训练结果：获得通用人体动作先验}

Kinetics-400 包含大量不同类型的人体动作，因此其主要作用并不是直接解决 Campus6 的六分类问题，而是让 ProtoGCN 提前获得较为稳定的人体运动先验。

经过 Kinetics-400 训练后，模型能够学习：

\begin{itemize}
\item 常见人体关节之间的空间关联；
\item 不同运动速度与运动幅度；
\item 上肢、下肢和躯干之间的动态协同；
\item 多种动作对应的典型 prototype 关系模式；
\item 具有类别判别能力的高层骨架表示。
\end{itemize}

在这一阶段，模型使用 Kinetics 的 400 类分类任务进行训练，最终保存训练得到的 ProtoGCN Backbone 权重：

\[
\theta_{\text{K400}}.
\]

该参数作为第二阶段 Campus6 模型的初始化：

\[
\theta_{\text{Campus6}}^{(0)}
\leftarrow
\theta_{\text{K400}}.
\]


\section{阶段二：Campus6 六分类全参数微调}

第二阶段使用校园行为数据对 Kinetics 预训练模型进行全参数微调。

首先加载 Kinetics-400 预训练得到的 ProtoGCN Backbone，然后将原来的：

\[
\text{-class Head}
\]

替换为：

\[
\text{-class Head}.
\]

当前六类目标行为包括：

\[
\{
\text{normal walk},
\text{normal run},
\text{playful chase},
\text{playful push},
\text{conflict chase},
\text{conflict push}
\}.
\]

这里采用的是 Full Fine-tuning，并不是只训练最后的六分类层。

也就是说，包括 GCN、MSTCN、动态拓扑学习模块以及 PRN 在内的 Backbone 参数都会继续根据 Campus6 数据进行反向传播更新。

当前 Campus6 工程正是从 Kinetics Skeleton 2D 的 ProtoGCN Backbone 初始化模型，并使用 6 类分类头继续训练；微调阶段采用较小学习率 0.0030.003，训练 40 epochs。

这种方式使模型首先从 Kinetics 学习“通用人体运动”，随后再将这些运动知识适配至更加细粒度的校园行为。


\section{GAP：利用动作自然语言描述辅助 Campus6 微调}

仅依赖骨架运动结构时，某些校园行为仍然存在明显的语义模糊问题。

例如：

\[
\text{playful chase}
\]

和

\[
\text{conflict chase}
\]

在骨架层面都可能表现为“两个人跑动”。

类似地：

\[
\text{playful push}
\]

和

\[
\text{conflict push}
\]

都可能出现“双人接触、上肢运动和身体位置变化”。

因此，仅告诉网络“这是哪一个类别”可能还不足以显式表达两个类别究竟差在哪里。

为此，本系统借鉴 Generative Action Description Prompts for Skeleton-based Action Recognition（GAP） 的思想，将自然语言动作描述作为辅助监督引入训练阶段。

GAP 的核心思想是将骨架特征与文本特征映射至同一个语义空间，通过特征相似度建立骨架动作与自然语言语义之间的对应关系。

对于两个已经 L2 归一化的特征：

\[
\hat{\mathbf x}_1,
\qquad
\hat{\mathbf x}_2,
\]

二者之间的内积即为余弦相似度：

\[
s
=
\hat{\mathbf x}_1
\hat{\mathbf x}_2^{\top}.
\]

原始 GAP 方法进一步利用标签关系构造 Ground Truth 相似度矩阵，并使用 KL 散度约束预测语义关系，其 KL 散度为：

\[
D_{\mathrm{KL}}(P\parallel Q)
=
\sum_{x\in\mathcal X}
P(x)
\log
\left(
\frac{P(x)}{Q(x)}
\right).
\]

上述思想与项目已有 GAP 学习页面中所总结的“特征归一化—相似度矩阵—语义对齐”过程保持一致。

\subsection{Campus6 身体部位级动作语义 Prompt}

为了向骨架特征提供更加明确的行为语义监督，本系统针对 Campus6 六类行为分别构造 Global、Head、Arms、Torso、Legs 五个身体层级的动作描述。与只使用类别名称相比，这些描述能够显式表达动作速度、运动幅度、身体协同关系以及双人相对距离变化等细粒度判别依据。

\begin{center}
\small
\begin{longtable}{|>{\raggedright\arraybackslash}p{0.15\linewidth}|>{\raggedright\arraybackslash}p{0.15\linewidth}|>{\raggedright\arraybackslash}p{0.15\linewidth}|>{\raggedright\arraybackslash}p{0.15\linewidth}|>{\raggedright\arraybackslash}p{0.15\linewidth}|>{\raggedright\arraybackslash}p{0.15\linewidth}|}
\hline
行为类别 & Global & Head & Arms & Torso & Legs \\
\hline
Normal Walk & a person walking & head moves with body direction & arms swing alternately forward and backward with low amplitude and small elbow bending & torso translates steadily forward & legs step alternately with short strides and small knee bending \\
\hline
Normal Run & a person running & head moves with body direction & arms swing quickly forward and backward with clear elbow bending & torso leans and translates quickly forward & legs stride quickly in alternation with clear knee bending \\
\hline
Playful Chase & two people running & no fixed motion pattern & no fixed motion pattern & no fixed motion pattern & no fixed motion pattern \\
\hline
Playful Push & two people make contact & heads remain near the interacting torsos & no fixed motion pattern & the two torsos do not move clearly farther apart & feet make small adjustments for balance \\
\hline
Conflict Chase & one person runs quickly behind another person & no fixed motion pattern & no fixed motion pattern & the distance between the two people changes quickly & no fixed motion pattern \\
\hline
Conflict Push & two people make contact and the distance between them changes substantially & no fixed motion pattern & no fixed motion pattern & after contact the distance between the two people changes substantially & no fixed motion pattern \\
\hline
\end{longtable}
\end{center}

这里的 Prompt 设计并不要求每一个身体部位都必须具有明确、固定的运动模式。对于缺少稳定局部规律的行为，直接使用 “no fixed motion pattern”，避免为了增加文本信息而人为引入并不存在的动作先验。

\section{文本描述的语义编码}

上述六类行为，每个类别包含 5 条身体部位描述，因此共形成：

\[
6\times5=30
\]

条动作语义描述。

使用冻结的 CLIP 文本编码器对这些描述进行编码，得到：

\[
\mathbf T
\in
\mathbb R^{6\times5\times512}.
\]

其中：

66

表示六个 Campus6 行为类别；

55

表示 Global、Head、Arms、Torso 和 Legs 五个身体语义区域；

512512

表示文本语义特征维度。

这些文本特征在训练前预先生成，并作为固定的语义监督目标保存，因此在 Campus6 训练过程中不需要继续更新 CLIP 文本编码器。

当前工程中的 GAP recognizer 实际加载的文本特征形状也正是：

[6,5,512].[6,5,512].


\section{身体部位骨架特征提取}

设 ProtoGCN 输出的高层骨架特征为：

\[
\mathbf F
\in
\mathbb R^{N\times M\times C\times T\times V}.
\]

其中：

\begin{itemize}
\item NN：batch size；
\item MM：人物数量；
\item CC：特征通道数；
\item TT：时间长度；
\item VV：关键点数量。
\end{itemize}

为了与五组文本语义进行对齐，将人体关键点划分为：

\[
P=
\{
P_{global},
P_{head},
P_{arms},
P_{torso},
P_{legs}
\}.
\]

对于某一个身体区域 pp，在人物维、时间维和该区域对应的关节维上执行平均池化：

\[
f_i^{(p)}
=
\operatorname{AvgPool}_{M,T,V_p}
(
\mathbf F_i
).
\]

这样即可获得当前样本在某个身体区域上的骨架动作表示。

需要注意的是，架构图中的“手腕”建议统一修改为 “上肢 Arms”，因为实际工程中的这一分支不仅包含 wrist，还包含 shoulder、elbow 和 wrist，即完整的上肢运动信息。


\section{骨架—文本语义对齐}

由于 ProtoGCN 骨架特征与 CLIP 文本特征位于不同的表示空间，因此对于每一个身体区域设置一个可学习线性投影：

\[
\mathbf v_i^{(p)}
=
\mathbf W_p
\mathbf f_i^{(p)}.
\]

将其投影至 512 维文本语义空间：

\[
v_i^{(p)}
\in
\mathbb R^{512}.
\]

随后对骨架特征和文本特征分别执行 L2 归一化：

\[
\hat{\mathbf v}_i^{(p)}
=
\frac{
\mathbf v_i^{(p)}
}{
\left\|
\mathbf v_i^{(p)}
\right\|_2
},
\]

\[
\hat{\mathbf t}_c^{(p)}
=
\frac{
\mathbf t_c^{(p)}
}{
\left\|
\mathbf t_c^{(p)}
\right\|_2
}.
\]

由此计算第 \(i\) 个骨架样本与第 \(c\) 类文本描述之间的相似度：

\[
\frac{
\hat{\mathbf v}_i^{(p)}
\cdot
\hat{\mathbf t}_c^{(p)}
}{
\tau
}.
\]

其中：

\[
\tau
\]

为温度参数。

当前项目设置：

\[
\tau=0.07.
\]

因此，对于一个属于 \(l_r\)un 的骨架样本，理想状态下其 Legs 特征应更加接近：

\begin{quote}
legs stride quickly in alternation with clear knee bending

\end{quote}

而不应该更加接近 normal walk 中：

\begin{quote}
legs step alternately with short strides and small knee bending

\end{quote}

这样模型就不仅知道两个样本分别属于哪个类别，还能够学习：

\[
\text{short strides}
\leftrightarrow
\text{fast strides},
\]

以及：

\[
\text{small knee bending}
\leftrightarrow
\text{clear knee bending}
\]

这样的更细粒度语义差异。


\section{GAP 身体部位语义损失}

对于每个身体区域，都可以得到当前骨架样本与六个文本类别之间的相似度：

\[
\mathbf s_i^{(p)}
=
[
s_{i,1}^{(p)},
s_{i,2}^{(p)},
\ldots,
s_{i,6}^{(p)}
].
\]

利用真实行为类别 \(y_i\)，分别计算各身体区域的语义分类损失：

\[
\mathcal L_{GAP}^{(p)}
=
CE
(
\mathbf s_i^{(p)},
y_i
).
\]

随后对五个身体区域求平均：

\[
\boxed{
\mathcal L_{GAP}
=
\frac{1}{5}
\sum_{p=1}^{5}
\mathcal L_{GAP}^{(p)}
}
\]

因此每一个训练样本实际上会同时收到五组语义约束：

\[
\text{Global}
+
\text{Head}
+
\text{Arms}
+
\text{Torso}
+
\text{Legs}.
\]

这种方式保留了 GAP 中“骨架特征—动作文本特征语义对齐”的核心思想，但针对 Campus6 已知六类标签的监督学习任务，采用 class-wise Cross Entropy 完成更加直接的语义约束。当前 RecognizerGCNGAP 的实际实现也是对 5 个身体区域分别计算相似度和 Cross Entropy，再取平均。

\section{Campus6 微调阶段联合损失}

Campus6 阶段最终同时受到三类监督。

第一部分为六分类 Cross Entropy：

\[
\mathcal L_{CE}.
\]

它负责保证模型能够正确预测最终行为类别。

第二部分为 ProtoGCN 的类别特异对比损失：

\[
\mathcal L_{CSC}.
\]

它利用 PRN 重构后的高层图关系约束同类动作更加聚合、不同动作更加分离。

第三部分为 GAP 身体部位动作语义损失：

\[
\mathcal L_{GAP}.
\]

它利用动作自然语言描述约束骨架特征具有更加明确的行为语义。

因此 Campus6 微调阶段的最终优化目标为：

\[
\boxed{
\mathcal L_{\text{Campus6}}
=
\mathcal L_{CE}
+
\lambda_{CSC}
\mathcal L_{CSC}
+
\lambda_{GAP}
\mathcal L_{GAP}
}
\]

当前工程中：

\[
C=0.2,\lambda_{CSC}=0.2,
\]

\[
P=0.15,\lambda_{GAP}=0.15,
\]

\[
\tau=0.07.
\]

因此：

\[
\boxed{
\mathcal L_{\text{Campus6}}
=
\mathcal L_{CE}
+
0.2\mathcal L_{CSC}
+
0.15\mathcal L_{GAP}
}
\]

这些参数与当前 Campus6 GAP 配置保持一致。

